\rhead{SM401: Technical Seminar}
\phantomsection
\section*{Technical Seminar (SM401)}
\label{course:SM401}

\noindent \small{\hyperref[main:scheme]{$\leftarrow$ Back to Scheme}} \vspace{0.5cm}

\noindent \textbf{Credits:} 2 (0-2-0)

\subsection*{Course Directive}
\noindent The Technical Seminar is a structured oral and written presentation course conducted in Semester 7. It is designed to develop the student's ability to communicate technical work rigorously and professionally to a peer and faculty audience. There is no fixed syllabus; the content is driven entirely by the student's own technical experience acquired during the preceding summer.

\vspace{0.3cm}

\noindent\textbf{Seminar Topic Selection} \\
Each student must present on \textit{one} of the following two categories:

\vspace{0.2cm}

\noindent\textbf{Category A — Industry or Research Internship:} Students who have completed a minimum 6-week internship (industry or research institution) during the summer following Year 3 must present the technical work undertaken during that internship. The presentation must go beyond a descriptive account and must include: the problem statement and its industrial or research context, the technical approach and tools employed, results and quantitative outcomes, and a critical reflection on limitations and future directions.

\vspace{0.3cm}

\noindent\textbf{Category B — Faculty-Supervised Minor Project:} Students who have not undertaken an external internship must complete a minor project under the supervision of a faculty member from their department. The project must involve a non-trivial technical deliverable — a working prototype, a simulation study, a literature survey with original analysis, or a dataset and its analysis — and must be initiated no later than the end of Semester 6. The seminar presentation reports the objectives, methodology, implementation, and findings of this project.

\newpage
\rhead{Curriculum Schema}
